\documentclass[CJK]{beamer}
\usepackage{CJKutf8}
\usepackage{beamerthemesplit}
\usetheme{Malmoe}
\useoutertheme[footline=authortitle]{miniframes}
\usepackage{amsmath}
\usepackage{amssymb}
\usepackage{graphicx}
\usepackage{eufrak}
\usepackage{color}
\usepackage{slashed}
\usepackage{simplewick}
\usepackage{tikz}
\usepackage{tcolorbox}
\graphicspath{{../figures/}}
%%figures
\def\lfig#1#2{\includegraphics[width=#1 in]{#2}}
\def\addfig#1#2{\begin{center}\includegraphics[width=#1 in]{#2}\end{center}}
\def\wulian{\includegraphics[width=0.18in]{emoji_wulian.jpg}}
\def\bigwulian{\includegraphics[width=0.35in]{emoji_wulian.jpg}}
\def\bye{\includegraphics[width=0.18in]{emoji_bye.jpg}}
\def\bigbye{\includegraphics[width=0.35in]{emoji_bye.jpg}}
\def\huaixiao{\includegraphics[width=0.18in]{emoji_huaixiao.jpg}}
\def\bighuaixiao{\includegraphics[width=0.35in]{emoji_huaixiao.jpg}}
%% colors
\def\blacktext#1{{\color{black}#1}}
\def\bluetext#1{{\color{blue}#1}}
\def\redtext#1{{\color{red}#1}}
\def\darkbluetext#1{{\color[rgb]{0,0.2,0.6}#1}}
\def\skybluetext#1{{\color[rgb]{0.2,0.7,1.}#1}}
\def\cyantext#1{{\color[rgb]{0.,0.5,0.5}#1}}
\def\greentext#1{{\color[rgb]{0,0.7,0.1}#1}}
\def\darkgray{\color[rgb]{0.2,0.2,0.2}}
\def\lightgray{\color[rgb]{0.6,0.6,0.6}}
\def\gray{\color[rgb]{0.4,0.4,0.4}}
\def\blue{\color{blue}}
\def\red{\color{red}}
\def\green{\color{green}}
\def\darkgreen{\color[rgb]{0,0.4,0.1}}
\def\darkblue{\color[rgb]{0,0.2,0.6}}
\def\skyblue{\color[rgb]{0.2,0.7,1.}}
%%control
\def\be{\begin{equation}}
\def\ee{\nonumber\end{equation}}
\def\bea{\begin{eqnarray}}
\def\eea{\nonumber\end{eqnarray}}
\def\bch{\begin{CJK}{UTF8}{gbsn}}
\def\ech{\end{CJK}}
\def\bitem{\begin{itemize}}
\def\eitem{\end{itemize}}
\def\bcenter{\begin{center}}
\def\ecenter{\end{center}}
\def\bex{\begin{minipage}{0.2\textwidth}\includegraphics[width=0.6in]{jugelizi.png}\end{minipage}\begin{minipage}{0.76\textwidth}}
\def\eex{\end{minipage}}
\def\chtitle#1{\frametitle{\bch#1\ech}}
\def\bmat#1{\left(\begin{array}{#1}}
\def\emat{\end{array}\right)}
\def\bcase#1{\left\{\begin{array}{#1}}
\def\ecase{\end{array}\right.}
\def\bmini#1{\begin{minipage}{#1\textwidth}}
\def\emini{\end{minipage}}
\def\tbox#1{\begin{tcolorbox}#1\end{tcolorbox}}
\def\pfrac#1#2#3{\left(\frac{\partial #1}{\partial #2}\right)_{#3}}
%%symbols
\def\sone{$\star$}
\def\stwo{$\star\star$}
\def\sthree{$\star\star\star$}
\def\sfour{$\star\star\star\star$}
\def\sfive{$\star\star\star\star\star$}
\def\rint{{\int_\leftrightarrow}}
\def\roint{{\oint_\leftrightarrow}}
\def\stdHf{{\textit{\r H}_f}}
\def\deltaH{{\Delta \textit{\r H}}}
\def\ii{{\dot{\imath}}}
\def\skipline{{\vskip0.1in}}
\def\skiplines{{\vskip0.2in}}
\def\lagr{{\mathcal{L}}}
\def\hamil{{\mathcal{H}}}
\def\vecv{{\mathbf{v}}}
\def\vecx{{\mathbf{x}}}
\def\vecy{{\mathbf{y}}}
\def\veck{{\mathbf{k}}}
\def\vecp{{\mathbf{p}}}
\def\vecn{{\mathbf{n}}}
\def\vecA{{\mathbf{A}}}
\def\vecP{{\mathbf{P}}}
\def\vecsigma{{\mathbf{\sigma}}}
\def\hatJn{{\hat{J_\vecn}}}
\def\hatJx{{\hat{J_x}}}
\def\hatJy{{\hat{J_y}}}
\def\hatJz{{\hat{J_z}}}
\def\hatj#1{\hat{J_{#1}}}
\def\hatphi{{\hat{\phi}}}
\def\hatq{{\hat{q}}}
\def\hatpi{{\hat{\pi}}}
\def\vel{\upsilon}
\def\Dint{{\mathcal{D}}}
\def\adag{{\hat{a}^\dagger}}
\def\bdag{{\hat{b}^\dagger}}
\def\cdag{{\hat{c}^\dagger}}
\def\ddag{{\hat{d}^\dagger}}
\def\hata{{\hat{a}}}
\def\hatb{{\hat{b}}}
\def\hatc{{\hat{c}}}
\def\hatd{{\hat{d}}}
\def\hatN{{\hat{N}}}
\def\hatH{{\hat{H}}}
\def\hatp{{\hat{p}}}
\def\Fup{{F^{\mu\nu}}}
\def\Fdown{{F_{\mu\nu}}}
\def\newl{\nonumber \\}
\def\vece{\mathrm{e}}
\def\calM{{\mathcal{M}}}
\def\calT{{\mathcal{T}}}
\def\calR{{\mathcal{R}}}
\def\barpsi{\bar{\psi}}
\def\baru{\bar{u}}
\def\barv{\bar{\upsilon}}
\def\qeq{\stackrel{?}{=}}
\def\torder#1{\mathcal{T}\left(#1\right)}
\def\rorder#1{\mathcal{R}\left(#1\right)}
\def\contr#1#2{\contraction{}{#1}{}{#2}#1#2}
\def\trof#1{\mathrm{Tr}\left(#1\right)}
\def\trace{\mathrm{Tr}}
\def\comm#1{\ \ \ \left(\mathrm{used}\ #1\right)}
\def\tcomm#1{\ \ \ (\text{#1})}
\def\slp{\slashed{p}}
\def\slk{\slashed{k}}
\def\calp{{\mathfrak{p}}}
\def\veccalp{\mathbf{\mathfrak{p}}}
\def\Tthree{T_{\tiny \textcircled{3}}}
\def\pthree{p_{\tiny \textcircled{3}}}
\def\dbar{{\,\mathchar'26\mkern-12mu d}}
\def\erf{\mathrm{erf}}
\def\const{\mathrm{constant}}
\def\pheat{\pfrac p{\ln T}V}
\def\vheat{\pfrac V{\ln T}p}
%%units
\def\fdeg{{^\circ \mathrm{F}}}
\def\cdeg{^\circ \mathrm{C}}
\def\atm{\,\mathrm{atm}}
\def\angstrom{\,\text{\AA}}
\def\SIL{\,\mathrm{L}}
\def\SIkm{\,\mathrm{km}}
\def\SIyr{\,\mathrm{yr}}
\def\SIGyr{\,\mathrm{Gyr}}
\def\SIV{\,\mathrm{V}}
\def\SImV{\,\mathrm{mV}}
\def\SIeV{\,\mathrm{eV}}
\def\SIkeV{\,\mathrm{keV}}
\def\SIMeV{\,\mathrm{MeV}}
\def\SIGeV{\,\mathrm{GeV}}
\def\SIcal{\,\mathrm{cal}}
\def\SIkcal{\,\mathrm{kcal}}
\def\SImol{\,\mathrm{mol}}
\def\SIN{\,\mathrm{N}}
\def\SIHz{\,\mathrm{Hz}}
\def\SIm{\,\mathrm{m}}
\def\SIcm{\,\mathrm{cm}}
\def\SIfm{\,\mathrm{fm}}
\def\SImm{\,\mathrm{mm}}
\def\SInm{\,\mathrm{nm}}
\def\SImum{\,\mathrm{\mu m}}
\def\SIJ{\,\mathrm{J}}
\def\SIW{\,\mathrm{W}}
\def\SIkJ{\,\mathrm{kJ}}
\def\SIs{\,\mathrm{s}}
\def\SIkg{\,\mathrm{kg}}
\def\SIg{\,\mathrm{g}}
\def\SIK{\,\mathrm{K}}
\def\SImmHg{\,\mathrm{mmHg}}
\def\SIPa{\,\mathrm{Pa}}

\def\courseurl{https://github.com/zqhuang/SYSU\_TD}

\def\tpage#1#2{
\begin{frame}
\begin{center}
\begin{Large}
\bch
热学 \\
第#1讲 #2

{\vskip 0.3in}

黄志琦

\ech
\end{Large}
\end{center}

\vskip 0.2in

\bch
教材：《热学》第二版，赵凯华，罗蔚茵，高等教育出版社
\ech

\bch
课件下载
\ech
\courseurl
\end{frame}
}

\title{Lesson 13 Conditions for Equilibrium }
  \author{}
  \date{}
\begin{document}
\tpage{13}{平衡判据}

\section{Review}

\begin{frame}
\chtitle{上讲内容回顾}
\bch 
\bitem
\item{热力学第二定律的四种表述}
\item{克劳修斯不等式和热力学第二定律的“微分”表示}
\item{熵的计算}
\eitem
\ech
\end{frame}

\begin{frame}
\chtitle{本讲内容}
\bch
\bitem
\item{眼花撩乱的平衡判据和背书宝的福音}
\item{真的懂了吗}
\eitem
\ech
\end{frame}


\section{A Rule of Thumb}

\begin{frame}
\chtitle{眼花撩乱的平衡判据}
\bch
教材从198页开始叨念各种（完全看不懂的）“平衡判据”：

\bitem
\item{熵判据：在内能和体积不变的情况下，对于一切可能的变动来说，平衡态的熵最大。}
\item{自由能判据：在温度和体积不变的情况下，对于一切可能的变动来说，平衡态的自由能最小。}
\item{自由焓判据：在温度和压强不变的情况下，对于一切可能的变动来说，平衡态的自由焓最小。}
\eitem

\addfig{1.5}{bzdsh.jpg}
\ech
\end{frame}

\begin{frame}
\chtitle{背书宝的福音}
\bch
\addfig{1}{beishubao.jpg}

要写出“$X$判据”，

\bitem
\item{第一步：把$X$写成全微分：$dX + \ldots dY + \ldots dZ = 0$}
\item{第二步：背书。在$Y$和$Z$不变的情况下，对于一切可能的变动来说，平衡态的$X$最$\ldots$。}
\item{第三步：判定是最大还是最小。$dX$, $dY$, $dZ$中必然有一个能量型的量（指$dU, dH, dF, dG$中的一个)，它前面系数为正则为最小，为负则为最大。}
\eitem

{\scriptsize
似乎有点理解了：

$Y$和$Z$不变，则$dX = 0$，代表$X$是某种极值？

能量型的量当然有下限？
}
\ech
\end{frame}


\begin{frame}
\chtitle{例子1：熵判据}
\bch
写出
$$dS - \frac{1}{T} dU - \frac{p}{T} dV = 0 $$
$dU$前系数为负，取最大。

在内能$U$和体积$V$不变的情况下，对于一切可能的变动来说，平衡态的熵$S$最大。

\skiplines

熵判据也常常说成，孤立系的平衡态熵最大。内能和体积不变，是孤立系的一种。
\ech
\end{frame}


\begin{frame}
\chtitle{例子2：自由能判据}
\bch
写出
$$dF  + S dT + p dV = 0 $$
$dF$前系数为正，取最小。


在温度$T$和体积$V$不变的情况下，对于一切可能的变动来说，平衡态的自由能$F$最小。

\ech
\end{frame}


\begin{frame}
\chtitle{例子3：自由焓判据}
\bch
写出
$$dG + S dT - V dp = 0$$
$dG$前系数为正，取最小。

在温度$T$和压强$p$不变的情况下，对于一切可能的变动来说，平衡态的自由焓$G$最小。

\ech
\end{frame}


\begin{frame}
\chtitle{例子4: 压强判据之一}
\bch
写出
$$dp - \frac{1}{V} dH + \frac{T}{V} dS $$
$dH$前系数为负，取最大。

在焓$H$和熵$S$不变的情况下，对于一切可能的变动来说，平衡态的压强$p$最大。

\ech
\end{frame}


\section{Understanding the Statements}

\begin{frame}
\chtitle{真懂了吗}
\bch
“一切可能的变动”是什么意思？

\addfig{1.5}{zhendongma.jpg}
\ech
\end{frame}

\begin{frame}
\chtitle{理解熵判据}
\bch
熵判据里“一切可能的变动”好理解：就是指孤立系的平衡态和一切非平衡态。对非平衡态，广延量内能和体积都有定义。


\addfig{1}{meimaobing.jpg}

证明：根据孤立系的熵增大原理，任何非平衡态的熵会持续增大直至到达平衡态。显然平衡态的熵是最大的。

\ech
\end{frame}


\begin{frame}
\chtitle{理解自由能判据}
\bch
自由能判据的条件是定温定体。“一切可能的变动”也包括非平衡态。这时需要假设系统和恒温的理想传导热库接触来理解非平衡态的温度。

\addfig{1}{meimaobing.jpg}

证明：从非平衡态到平衡态的过程为不可逆过程。定体过程不对外做功，又由于恒温条件下对外做功小于自由能的减少（对不可逆过程）。故自由能的减少大于零。即非平衡态的自由能大于平衡态的自由能。

\ech
\end{frame}

\begin{frame}
\chtitle{理解自由焓判据}
\bch
自由焓判据的条件是定温定压。“一切可能的变动”也包括非平衡态。这时需要假设系统和恒温的理想传导热库接触来理解非平衡态的温度。

\addfig{1}{meimaobing.jpg}

证明：从非平衡态到平衡态的过程为不可逆过程。定温定压条件下，对不可逆过程，其他形式的做功小于自由焓的减少，又这里不考虑其他形式的做功（即没有额外安装发电机什么的），故自由焓的减少大于零。即非平衡态的自由焓大于平衡态的自由焓。

\ech
\end{frame}

\begin{frame}
\chtitle{三大平衡条件}
\bch
\bitem
\item[1]{\blue 热平衡条件：系统内部温度均匀}
\item[2]{\blue 力学平衡条件：系统内部压强均匀}
\item[3]{\blue 相平衡条件：系统内各相化学势相等}
\eitem

{\scriptsize
注：第三条仅考虑了物态变化（如液态水和水蒸气之间的平衡），其实它可以推广到涉及化学反应的一般的化学平衡条件。化学反应比相变更为复杂，各种反应物或者产物有可能混合在一起产生额外的混合熵（请参考教材208页）。
}
\ech
\end{frame}


\begin{frame}
\chtitle{热平衡条件}
\bch
热平衡条件可以从孤立系的熵判据推出：如果有相邻的温度不相等的子系统，则由高温到低温的热量传输必然增大熵。这说明系统还未达到平衡态（熵最大的状态）。
\addfig{1.5}{eqc1.jpg}
$$\Delta S = \frac{-Q}{T_1} + \frac{Q}{T_2} > 0$$
\ech
\end{frame}

\begin{frame}
\chtitle{力学平衡条件}
\bch
力学平衡条件可以从定温定体系统的自由能判据推出：如果有相邻的压强不相等的子系统，则压强大的子系统等温膨胀$\Delta V$，压强小的子系统等温收缩$\Delta V$，可使系统温度和体积不变的情况下总自由能减小。这说明系统还未达到平衡态（自由能最低的态）。
\addfig{1.5}{eqc2.jpg}
$$\Delta F = -p_1 \Delta V + p_2 \Delta V < 0 $$
\ech
\end{frame}


\begin{frame}
\chtitle{相平衡条件}
\bch
设有种可以相互转化的物态。每种物态的粒子数分别为$N_1$, $N_2$, $\ldots$, $N_n$。每种物态成分的化学势即该成分的平均每个粒子的自由焓，即
$$ \mu_i \equiv \pfrac G{N_i}{T,p,N_1, N_2, \ldots, N_{i-1}, N_{i+1}, \ldots, N_n} $$

\addfig{1.5}{eqc3.jpg}

相平衡条件可以从定温定压系统的自由焓判据推出：如果定温定压条件下$ \mu_i > \mu_j$，则可以通过$i\rightarrow j$的相变降低系统的总自由焓。
，这说明系统还未达到平衡态。
\ech
\end{frame}


\begin{frame}
\chtitle{经典理想气体的化学势}
\bch
{\small
考虑经典的理想气体，在$i$态上的分子数为
$$ n_i = e^{-\frac{\varepsilon_i-\mu}{kT}}$$
下面我们证明，这里的$\mu$和用自由焓定义的化学势相同}

{\scriptsize
  我们来计算温度为$T$的处于平衡态的单一成分理想气体的自由焓。
  $$G = U + pV - TS$$
  设第i个态上有$n_i$个分子，总分子数$N=\sum_i n_i$。因为要用到熵的绝对数值，我们就不能像以前那样直接忽略粒子的全同性了（混合气体的熵不同于单一成分气体的熵）。因为置换$N$个粒子共有$N!$种可能性。按照熵是状态数的对数的理解，按可区分粒子计算出来的熵则要减去$\ln (N!) \approx N \ln N - N $。最后，按照热力学的单位制，熵以$k^{-1}$为单位，于是
  $$\frac{S}{k} = N\sum_i\left(-\frac{n_i}{N}\ln\frac{n_i}{N}\right)  - (N\ln N - N) = N-\sum_i (n_i\ln n_i) = N - \sum_in_i\frac{\mu - \varepsilon_i}{kT} =N+\frac{U-\mu N}{kT}  $$
  即
  $$TS = kNT + U - \mu N = pV + U - \mu N$$
  移项即得结论$$ G = \mu N$$
}
\ech
\end{frame}

\begin{frame}
\chtitle{光子气体的化学势}
\bch
利用光子气体
$$U = aV T^4$$
以及熵
$$ S = \frac{4}{3}aV T^3$$
压强
$$ p = \frac{a}{3} T^4$$
很容易得到自由焓
$$ G = U + pV - TS = 0$$
即化学势为零。

但，光子气体的$U = aVT^4$本身就是我们用化学势为零的假设推导出来的，为了避免循环论证，我们必须用不同的方法进行证明。
\ech
\end{frame}


\begin{frame}
\chtitle{光子气体内能公式的独立推导}
\bch
{\small
把光子气体理解为温度为$T$的黑体辐射，则光子气体的能量密度只能是$T$的函数。我们假设
$$ U = V f(T)$$
其中$f$为待定函数。
我们曾经推导了极端相对论气体的压强公式，并得到了$p = \frac{1}{3} n\overline{\varepsilon} = \frac{U}{3V} $，即
$$p = \frac{1}{3} f(T)$$
由内能和状态方程的关系得到
$$ f(T) = \pfrac UVT = T\pfrac pTV - p = \frac{1}{3} \left[T f'(T) - f(T)\right]$$
化简即
$$ \frac{d \ln f}{d\ln T}  = 4 $$
即$ f = a T^4 $。这个证明无法确定积分常数$a$。但由此可以推出光子气体化学势为零，再按第5讲那样用玻色-爱因斯坦分布可以计算出$a$。
}
\ech
\end{frame}

\begin{frame}
  \chtitle{后记：关于不考的一些内容}
  \bch
  问题1: 教材204页至208页在瞎扯什么，完全看不懂……

  \skiplines
  
  答：是这样的，教材在这节开头忘写了这么一段：

  {\small 假设气态和液态都能用范德瓦尔斯方程描述（看成一种假想的单相态），我们可以计算出理论上的单相$F$-$V$曲线等。实际发生的情况未必和这种假想的单相态理论预言相同，下面我们用自由能判据等来进行判断……}

  \ech
\end{frame}


\begin{frame}
  \chtitle{后记：关于不考的一些内容}
  \bch
  问题2: 直接跳过第五章的知识，是否会对学习后续课程造成影响。

  \skiplines
  
  答：本科阶段肯定不会。研究生阶段基本不会。研究生之后的阶段——你想太多了\bye

  \ech
\end{frame}

\begin{frame}
  \chtitle{知识结构总结}
\addfig{4.3}{summary.png}  
\end{frame}


\section{Homework}

\begin{frame}
  \chtitle{第13周作业(序号接第12周)}
  \bch
  {\small
  \bitem
\item[33]{补充完整下列平衡判据：
  \bitem
\item[(1)]{在\_\_和\_\_不变的情况下，平衡态的焓$H$最\_\_。}
\item[(2)]{在自由能$F$和\_\_不变的情况下，平衡态的温度$T$最\_\_。}
\item[(3)]{在焓$H$和\_\_不变的情况下，平衡态的熵$S$最\_\_。}
\item[(4)]{在自由焓$G$和\_\_不变的情况下，平衡态的压强$p$最\_\_。}
  \eitem}
\item[34]{无重力环境下，定温定体容器中装满水蒸气，水蒸气中凝结出一个半径为$R$的悬浮球状液态水滴。这个系统可以看作气相，液相，气液交界处的表面相三相共存。表面相的方程为
  $$ dU = TdS + \sigma dA $$
其中$\sigma$为水的表面张力系数，$A = 4\pi R^2$为水滴表面积。利用自由能判据推导平衡态下液态水和气态水的压强之差(用$\sigma$, $R$表示)。}
\item[35]{上题中的相平衡条件还是液态水和气态水的化学势相等吗？为什么？}
  \eitem
  }
  \ech
\end{frame}

\end{document}
